\documentclass{article}
\usepackage[margin=1in]{geometry}
\usepackage{listings}
\usepackage{color}
\usepackage{hyperref}
\usepackage{enumitem}

\definecolor{codegreen}{rgb}{0,0.6,0}
\definecolor{codegray}{rgb}{0.5,0.5,0.5}
\definecolor{codepurple}{rgb}{0.58,0,0.82}
\definecolor{backcolour}{rgb}{0.95,0.95,0.92}

\lstdefinestyle{mystyle}{
	backgroundcolor=\color{backcolour},   
	commentstyle=\color{codegreen},
	keywordstyle=\color{magenta},
	numberstyle=\tiny\color{codegray},
	stringstyle=\color{codepurple},
	basicstyle=\ttfamily\footnotesize,
	breakatwhitespace=false,         
	breaklines=true,                 
	captionpos=b,                    
	keepspaces=true,                 
	numbers=left,                    
	numbersep=5pt,                  
	showspaces=false,                
	showstringspaces=false,
	showtabs=false,                  
	tabsize=2
}

\lstset{style=mystyle}

\title{ICSE Class 10 Python \\ Lists and Tuples - Exam Ready Notes}
\author{}
\date{}

\begin{document}
	
	\maketitle
	
	\section{Lists in Python}
	
	\begin{itemize}
		\item A \textbf{list} is a \textbf{mutable} (changeable), \textbf{ordered} sequence of elements.
		\item Created using \textbf{square brackets} \texttt{[]}.
		\item Can store elements of \textbf{different data types} (heterogeneous).
		\item Supports \textbf{indexing} (starts at 0) and \textbf{negative indexing}.
	\end{itemize}
	
	\subsection{Examples}
	\begin{lstlisting}[language=Python]
		# Empty list
		l1 = []
		
		# List with different types
		l2 = [10, 20.5, "Hello", True, [1,2,3]]
		
		# Accessing elements
		print(l2[0])     # 10
		print(l2[-1])    # [1, 2, 3]
		print(l2[2:4])   # ['Hello', True]  # Slicing
	\end{lstlisting}
	
	\subsection{List Operations}
	\begin{itemize}
		\item Concatenation: \texttt{list1 + list2}
		\item Repetition: \texttt{list * n}
		\item Membership: \texttt{element in list}
		\item Length: \texttt{len(list)}
	\end{itemize}
	
	\subsection{Important List Methods}
	
	\begin{tabular}{|l|l|l|}
		\hline
		\textbf{Method} & \textbf{Description} & \textbf{Example} \\
		\hline
		\texttt{append(x)} & Adds element at the end & \texttt{lst.append(50)} \\
		\texttt{insert(i, x)} & Inserts at index i & \texttt{lst.insert(1, 25)} \\
		\texttt{extend(lst2)} & Adds all elements of lst2 & \texttt{lst.extend([60,70])} \\
		\texttt{remove(x)} & Removes first occurrence of x & \texttt{lst.remove(20)} \\
		\texttt{pop(i)} & Removes and returns element at i & \texttt{lst.pop()} \\
		\texttt{clear()} & Removes all elements & \texttt{lst.clear()} \\
		\texttt{index(x)} & Returns index of first occurrence & \texttt{lst.index(30)} \\
		\texttt{count(x)} & Counts occurrences of x & \texttt{lst.count(10)} \\
		\texttt{sort()} & Sorts in ascending order & \texttt{lst.sort()} \\
		\texttt{reverse()} & Reverses the list & \texttt{lst.reverse()} \\
		\hline
	\end{tabular}
	
	\subsection{Nested Lists}
	\begin{lstlisting}[language=Python]
		nested = [[1,2,3], [4,5,6], [7,8,9]]
		print(nested[1][2])  # 6
	\end{lstlisting}
	
	\section{Tuples in Python}
	
	\begin{itemize}
		\item A \textbf{tuple} is an \textbf{immutable} (cannot be changed), \textbf{ordered} sequence.
		\item Created using \textbf{parentheses} \texttt{()} or comma separated values.
		\item Faster than lists; used for fixed data.
	\end{itemize}
	
	\subsection{Examples}
	\begin{lstlisting}[language=Python]
		t1 = (10, 20, 30)
		t2 = 1, 2, 3, 4          # Parentheses optional
		t3 = ("apple",)          # Single element needs comma
		empty = ()               # Empty tuple
	\end{lstlisting}
	
	\subsection{Tuple Assignment / Packing-Unpacking}
	\begin{lstlisting}[language=Python]
		# Packing
		t = 10, 20, 30
		
		# Unpacking
		a, b, c = t
		print(a, b, c)  # 10 20 30
		
		# Swap variables
		x, y = y, x
	\end{lstlisting}
	
	\section{Key Differences: List vs Tuple}
	
	\begin{tabular}{|l|l|l|}
		\hline
		\textbf{Feature} & \textbf{List} & \textbf{Tuple} \\
		\hline
		Mutability & Mutable & Immutable \\
		Syntax & \texttt{[]} & \texttt{()} \\
		Speed & Slower & Faster \\
		Use case & Changing data & Fixed data \\
		Memory & More & Less \\
		\hline
	\end{tabular}
	
	\section{Common Functions}
	\begin{itemize}
		\item \texttt{len()} $\rightarrow$ length
		\item \texttt{max()}, \texttt{min()} $\rightarrow$ on numeric sequences
		\item \texttt{sum()} $\rightarrow$ for numbers
		\item \texttt{sorted(seq)} $\rightarrow$ returns new sorted list
		\item \texttt{list(tuple)} / \texttt{tuple(list)} $\rightarrow$ type conversion
	\end{itemize}
	
\end{document}